\problemname{Grinding}
\noindent

Efter att ha testat på AI-verktyget ''PO kodbankare'' insåg Johan att jobbsäkerheten för programmerare är högst tveksam.
Därmed har han bytt karriär till äventyrare.
En äventyrares främsta arbetsuppgift är att döda monster.
Det finns $N$ stycken grottor, numrerade från $1$ till $N$.
Varje grotta har ett visst antal monster som står på rad.
Varje monster har en viss styrka och XP-värde.

Johan börjar med styrka $1$ och har endast tillgång till grotta $1$ i början.
För att döda monster kan han \texttt{anfalla} en grotta.
När han anfaller strider han först mot det första monstret.
Om han har samma eller mer styrka än monstret vinner han, och då ökar hans styrka med monstrets XP, och han går sedan vidare till nästa monster.
Om han någonsin förlorar en strid eller besegrar alla monster lämnar han grottan.
När han lämnar grottan kommer alla monster tillbaka, precis som de var från början.
Om han besegrar alla monster i grotta $i$ och $i \neq N$, så låser han upp grotta $i+1$, och kan anfalla den i framtiden.

Hans slutgiltiga mål är att kunna besegra demonkungen, som han vet kräver styrka som minst $B$.
Avgör minsta antalet anfall som han behöver göra för att få styrka som minst $B$.

\section*{Indata}
Den första raden av indata innehåller heltalen $N$ och $B$ ($1 \leq N \leq 2 \cdot 10^5$, $1 \leq B \leq 10^9$),
antalet grottor och din målstyrka.

Därefter följer beskrivningar av de $N$ grottorna.
Beskrivningen av varje grotta börjar med en rad som innehåller heltalet $M$ ($1 \leq M \leq 3 \cdot 10^5$), antalet monster i grottan.
Därefter följer $M$ rader som vardera innehåller heltalen $s$ och $e$ ($1 \leq s,e \leq 10^9$),
styrkan och XP-värdet av ett monster.
Dessa är angivna i samma ordning som du kommer strida mot monstrena i grottan.

Det är garanterat att första monstret i första grotten har styrka $1$.
Därmed går det alltid att uppnå styrka $B$.

Låt $T$ vara summan av alla $M$, det vill säga totala antalet monster. Då håller det att $1 \leq T \leq 3 \cdot 10^5$.

\section*{Utdata}
Skriv ut ett heltal: det minsta antalet anfall som Johan behöver göra för att styrka $B$ eller mer.

\section*{Poängsättning}
Din lösning kommer att testas på en mängd testfallsgrupper.
För att få poäng för en grupp så måste du klara alla testfall i gruppen.

\noindent
\begin{tabular}{| l | l | p{12cm} |}
  \hline
  \textbf{Grupp} & \textbf{Poäng} & \textbf{Gränser} \\ \hline
  $1$    & $10$       & $N, B, T \leq 100$ \\ \hline
  $2$    & $18$       & $N \leq 2000$, $B, T \leq 5000$ \\ \hline
  $3$    & $20$       & $B \leq 1000$ \\ \hline
  $4$    & $15$       & $N=1$ \\ \hline
  $5$    & $10$       & Varje grotta har som mest 10 monster i sig. \\ \hline
  $6$    & $27$       & Inga ytterligare begränsningar. \\ \hline
\end{tabular}

\section*{Förklaring av exempelfall $1$}
I exempelfall $1$ kan Johan anfalla grotta $1$.
Eftersom han börjar med styrka $1$ vinner han isåfall striden och ökar sin styrka med $1$.
Det inte finns fler grottor eller och därmed får Johan aldrig mer än $1$ styrka per anfall.
Alltså tar det $99$ stycken anfall att få $B=100$ eller mer styrka.

\section*{Förklaring av exempelfall $2$}
Till att börja med har Johan endast tillgång till grotta $1$.
Om han anfaller den besegrar han både monstrena och ökar sin styrka med $3$.
Han låser dessutom upp grotta $2$, men kan ännu inte besegra monstret i den.
Om han anfaller grotta $1$ två gånger till blir hans styrka $10$.
Han kan nu anfalla grotta $2$ två gånger till för att få $18$ styrka, vilket är tillräckligt.
Det går inte att utföra färre anfall.
